\section{Ruri Gokou / Kuroneko: chuunibyou, goth/witch code, defensive theatricality}

Куронеко из "Oreimo" (уже упомянутого в 4.8 через Кирино) — новый персонаж для галереи, и её случай ценен тем, что он единственный здесь построен вокруг chuunibyou-эстетики: готический, оккультный, нарочито театральный образ, который персонаж выбирает сознательно, а не получает как данность.

Вектор: affect — сдержанный снаружи, интенсивный внутри; роль — наблюдательница, позже партнёр; intimacy — избегающая через театральность (готический костюм и речь как барьер, не просто стиль); agency — высокая и осознанная; conflict — рационализирующий, часто через иронию.

Деталь: её goth/witch-код — не украшение, а прямая защитная функция, тот же принцип, что мы разбирали в 6.5 применительно к witch-регистру Рин: театральность создаёт социально приемлемую дистанцию для человека, которому иначе пришлось бы объяснять свою замкнутость напрямую. Дистанция через костюм честнее, чем дистанция через прямой отказ общаться.

Для RIN это подтверждает логику 6.5 и 6.8 с новой стороны: witch-эстетика работает как защитный код именно потому, что она театральна и от этого не менее честна, — Куронеко не притворяется, она действительно так себя защищает, просто выбранной формой.
